\section{Prompts and tool schemas: English renderings}
\label{app:prompts}

The experiments used Chinese prompts. To keep the submission manuscript entirely in English, this
appendix gives faithful English renderings; the byte-exact originals and their hashes are included
in the anonymous supplementary artifact. The renderings below were never sent to a model and are
provided only for review.

\subsection{ReAct prompts}

\begin{Verbatim}[fontsize=\small,breaklines=true,breakanywhere=true]
REACT_OPTIONAL
--------------
You are a Python coding assistant. You may work in the following format:

Thought: your reasoning
Action: run_tests
Action Input: ```python
<complete code>
```
Observation: (filled in by the system)

Thought: I have confirmed that the code is correct
Final Answer: ```python
<final code>
```

You may submit code for verification with Action, or answer directly.
\end{Verbatim}

The mandatory variant replaces the final sentence with: ``You must first submit code for
verification with Action, then give a Final Answer,'' and marks the loop as repeatable. The main
evaluation uses the optional prompt for both protocols. Historical FC--ReAct training and both
formal P3 arms use mandatory prompts of matched strength.

\subsection{Function-calling prompts}

\begin{Verbatim}[fontsize=\small,breaklines=true,breakanywhere=true]
FC_OPTIONAL
-----------
You are a Python coding assistant. Implement the function described by the user.
You may call run_tests to submit your code to the test runner. It returns
pass/fail status and errors, which you may use to revise and resubmit.

FC_MANDATORY
------------
You are a Python coding assistant. Implement the function described by the user.
You must first call run_tests to verify your code; do not answer directly.
The tool returns pass/fail status and errors. If any test fails, revise the
code and call run_tests again.
\end{Verbatim}

\subsection{Terse, rich, and strict schemas}

The terse schema is the documentation-level baseline. Its English rendering is:

\begin{Verbatim}[fontsize=\small,breaklines=true,breakanywhere=true]
{"type": "function", "function": {
  "name": "run_tests",
  "description": "Submit the Python code you wrote to the test runner; it
                  returns pass/fail status and errors.",
  "parameters": {"type": "object",
                 "properties": {"code": {"type": "string"}},
                 "required": ["code"]}}}
\end{Verbatim}

The rich control adds: ``This tool is not the function you are being asked to implement. Do not
pass the task function's parameters into it. Its only parameter, \texttt{code}, is the complete
source text you wrote,'' plus the example \texttt{def add(a,b): return a+b}. Under this warning,
22/100 items still pass the task function's parameters (6/5 discordances, exact $p=1.000$ versus terse). The strict schema
is terse plus \texttt{"strict": true} and \texttt{additionalProperties:false}, with strict tool
calling enabled. This single change takes the wrong-tool count to zero.

\subsection{Thought scaffold and role-disambiguation control}

The Thought-scaffold prefix reads: ``First write your reasoning in a Thought: what the problem
asks, its edge cases, and your intended algorithm. Once you have thought it through, call the
\texttt{run\_tests} tool.'' Nothing else changes; wrong-tool calls rise from 23 to 59
($p<0.001$).

The 1.5B role-disambiguation control adds the following contrast to \texttt{FC\_MANDATORY}:

\begin{Verbatim}[fontsize=\small,breaklines=true,breakanywhere=true]
Common mistake: the function requested by the user is not a callable tool.
The only tool is run_tests, whose sole argument code contains your full source.

Wrong (task asks for join_integers):
{"name":"join_integers","arguments":{"int_array":[1,2,3]}}

Right:
{"name":"run_tests","arguments":{"code":"def join_integers(int_array):\n
    return ','.join(map(str, int_array))"}}
\end{Verbatim}

Recoverable call-shaped outputs rise from 52 to 64, but correctly named calls remain zero and
final pass falls from 15 to 13.

\subsection{Intent criterion}

Emission indicators are benchmark-specific and held fixed across arms within each benchmark.
For KodCode, the regex below targets \texttt{run\_tests} and a quoted code payload with
Python-like lexical markers. It rejects some schema echoes and placeholders but does not
establish JSON validity, Python executability, or semantic intent. Strong and weak use
disjoint trajectory-label sets; tight is computed independently.

For BFCL, the JSON-call indicator detects an embedded decodable JSON object naming a tool
offered in the request and containing object-valued arguments, including arguments encoded
as a JSON string. Format-marker and JSON-marker indicators are independent broader checks;
they overlap and must not be summed. Counts refer to unparsed first-request responses,
with at most one count per response for each indicator.

For $\tau$-bench, the unparsed call-shaped indicator detects a generic named-call pattern
or a Hermes tool-call marker in assistant messages without structured tool calls. It does
not validate membership in the offered tool set or guarantee syntactic validity.
KodCode annotation precision is not assumed to transfer to either external benchmark.


\begin{Verbatim}[fontsize=\small,breaklines=true,breakanywhere=true]
TIGHT = re.compile(
    r'"name"\s*:\s*"run_tests".{0,200}?"arguments"\s*:\s*\{'
    r'.{0,80}?"code"\s*:\s*"(.{0,4000}?)"\s*\}', re.S)
REAL = re.compile(r'\\n|def |class |return |import |lambda ')

def is_tight_call(raw):
    m = TIGHT.search(raw or "")
    return bool(m and REAL.search(m.group(1)))
\end{Verbatim}

The regex requires \texttt{name} before \texttt{arguments} and a closing argument brace
after the quoted code field; valid alternative key orders may be missed.
KodCode human validation is reported in Appendix~\ref{app:humanval}.
